\documentclass[12pt,a4paper]{article}
\usepackage[top=23mm,bottom=24mm,left=27mm,right=27mm,headheight=15pt]{geometry}
\usepackage{amsmath,amssymb,booktabs,array,xcolor,fancyhdr}
\definecolor{ink}{HTML}{194D61}
\usepackage[hypertexnames=false,colorlinks=true,linkcolor=ink,urlcolor=ink,bookmarksnumbered=true]{hyperref}
\hypersetup{pdftitle={Kernel approximate dynamic programming: results and proofs},pdfauthor={Kernel ADP formalization project}}
\setlength{\parindent}{0pt}\setlength{\parskip}{5pt plus 1pt}
\setcounter{tocdepth}{1}
\raggedbottom
\pagestyle{fancy}\fancyhf{}\fancyhead[L]{\small\textcolor{ink}{Kernel approximate dynamic programming}}
\fancyhead[R]{\small\textcolor{ink}{\nouppercase{\leftmark}}}
\fancyfoot[C]{\thepage}\renewcommand{\headrulewidth}{0.3pt}
\renewcommand{\sectionmark}[1]{\markboth{#1}{}}
\newcommand{\partsection}[2]{\clearpage\section[#1]{#2}}
\newcommand{\lead}[1]{\par\smallskip\textbf{#1}\enspace}
\newcommand{\E}{\mathbb E}\newcommand{\Prob}{\mathbb P}
\newcommand{\J}{\mathcal J_{C,B}}\newcommand{\Cbar}{\overline C}
\newcommand{\occ}{\pi_{\mu^*,\nu}}
\newcommand{\leanfile}[2]{\href{https://github.com/nikhilpb/kadp/blob/main/lean/#1.lean}{#2}}
\begin{document}
\thispagestyle{plain}
{\huge\bfseries Kernel approximate\\dynamic programming\par}
\medskip
{\Large Results and proofs for discounted cost minimization\par}
\bigskip
A kernel represents a value function through inner products of state
features. This lets us fit the function by solving a finite quadratic
program, even when the features live in an infinite-dimensional space.
There are two mathematical questions. Which value function does the
quadratic program recover? How well can a value function fitted on a
sample approximate the optimal cost?

The first question has an exact answer for each realized sample.
Proposition~1 recovers the fitted weight from any optimal multiplier,
and Proposition~2 characterizes optimality by transfers between pairs
of multipliers. The second question concerns an idealized sampled
program with bounded weights and intercept and a prescribed state law.
Theorems~1 and~2 bound its error by the best error available within
that class, plus a sampling term. Theorem~3 supplies the uniform
probability estimate that makes this possible.

\tableofcontents
\medskip
The numbering of propositions, theorems, lemmas and parenthesized
manuscript equations follows Bhat, Farias, Moallemi and Zheng,
\emph{Non-parametric Approximate Dynamic Programming via the Kernel
Method}, August 29, 2022. The final section locates the corrections
and identifies which guarantees they affect.

\partsection{Bellman LP}{The Bellman linear program}
\label{sec:bellman}
Let \(\mathcal S,\mathcal A\) be finite nonempty sets of states and actions.
An action \(a\) in state \(x\) incurs cost \(g_{x,a}\) and moves to \(y\) with
probability \(p(x,y,a)\). Each transition row has nonnegative entries summing
to one. Write \(P_aJ(x)=\sum_y p(x,y,a)J(y)\). The discount is \(0<\alpha<1\).
For a stationary policy \(\mu:\mathcal S\to\mathcal A\), the expected discounted
cost starting at \(x\) is \(J_\mu(x)=\mathbb E_{x,\mu}\sum_{t\ge0}\alpha^t
g_{X_t,\mu(X_t)}\). Put \(g_\mu(x)=g_{x,\mu(x)}\), \(P_\mu(x,y)=p(x,y,\mu(x))\), and set
\[
 (T_\mu J)(x)=g_\mu(x)+\alpha(P_\mu J)(x),\qquad
 (TJ)(x)=\min_{a\in\mathcal A}\{g_{x,a}+\alpha P_aJ(x)\}.
\]
Both operators preserve pointwise order. Stochastic averaging and taking a
finite minimum each preserve a uniform bound on differences, so both are
contractions in the sup norm with factor \(\alpha\).

\lead{Bellman LP and greedy policies.}
The optimal cost \(J^*(x)=\min_\mu J_\mu(x)\) is attained by one policy
simultaneously at all states. A policy attains it exactly when each of its
actions attains the minimum defining \(TJ^*\). For every vector of strictly
positive state weights \(\nu\), \(J^*\) is the unique solution of
\[
 \max_J\ \sum_x\nu_xJ(x)
 \qquad\text{subject to}\qquad
 J(x)\le g_{x,a}+\alpha P_aJ(x)\quad\text{for all }x,a. \tag{1}
\]

\lead{Proof.}
Contraction gives a unique fixed point \(U\) of \(T\). Choose a minimizing action at
each state to obtain \(\mu\) with \(T_\mu U=U\). Conditioning on the first
transition gives \(J_\mu=T_\mu J_\mu\), and uniqueness gives \(U=J_\mu\).
For every policy \(\sigma\), \(U=TU\le T_\sigma U\). Iteration and contraction
give \(U\le J_\sigma\), proving \(U=J^*\). Conversely, \(J_\sigma=J^*\)
implies \(T_\sigma J^*=TJ^*\), so every action selected by \(\sigma\) is greedy.

The constraints of (1) say \(J\le TJ\). Iteration gives \(J\le T^nJ\to J^*\)
componentwise. Since \(J^*=TJ^*\) is feasible, it maximizes the objective.
Any strict inequality \(J(x)<J^*(x)\) strictly lowers that objective because
\(\nu_x>0\). This proves uniqueness.

\lead{Why the two optimization signs differ.}
The policy minimizes cost. The LP maximizes a lower bound on that cost.
For one state, two self-looping actions with costs \(0,1\) and
\(\alpha=1/2\), the two policies cost \(0\) and \(2\). The LP constraints
are \(J\le J/2\) and \(J\le 1+J/2\), so its maximum is \(0\).
The policy maximization printed in \S2.1 would select the action costing \(2\).

\partsection{Sampled duality}{The sampled primal and its dual}
\label{sec:dual}
A feature map \(\Phi:\mathcal S\to\mathcal H\) takes values in a real Hilbert
space. Its inner product and norm are \(\langle\cdot,\cdot\rangle\) and
\(\|\cdot\|\). Fit \(J_{\mathbf z,b}(x)=\langle\Phi(x),\mathbf z\rangle+b\).
Let \(X_1,\ldots,X_N\) be \(N\ge1\) draws from a distribution \(\pi\) positive
on every state, and let \(w_i=\nu_{X_i}/\pi_{X_i}\) for positive relevance
weights \(\nu\). Fix the sample and give each repeat its own constraint row.
Write \(g_{ia}=g_{X_i,a}\) and define
\[
 m=N^{-1}\sum_i w_i\Phi(X_i),\qquad d_{ia}=\Phi(X_i)-\alpha P_a\Phi(X_i),
 \qquad v(\lambda)=\sum_{i,a}\lambda_{ia}d_{ia}.
\]
Here \(P_a\Phi(x)=\sum_y p(x,y,a)\Phi(y)\). For \(\Gamma,\kappa>0\) the
sampled primal program is
\[
 \begin{aligned}
 \max_{\mathbf z,b,s}\quad&P(\mathbf z,b,s)=\langle m,\mathbf z\rangle+b-
       \tfrac{\kappa}{N}\sum_i s_i-\tfrac{\Gamma}{2}\|\mathbf z\|^2\\
 \text{subject to}\quad&s_i\ge0,\qquad
    \langle d_{ia},\mathbf z\rangle+(1-\alpha)b\le g_{ia}+s_i
       \quad(\text{all }i,a).
 \end{aligned}\tag{6}
\]
The intercept has coefficient one. A slack \(s_i\) covers the largest
Bellman constraint violation among actions at draw \(i\).
Define the feasible multiplier set and the normalized dual objective by
\[
 \begin{aligned}
 \Lambda&=\{\lambda\ge0:\ \sum_a\lambda_{ia}\le\kappa/N,\
               \sum_{i,a}\lambda_{ia}=1/(1-\alpha)\},\\
 G(\lambda)&=\sum_{i,a}g_{ia}\lambda_{ia}+
               \tfrac1{2\Gamma}\|m-v(\lambda)\|^2 .
 \end{aligned}
\]

\lead{Proposition 1 (strong duality and recovery).}
If \(\kappa\ge1/(1-\alpha)\), the maximum of (6) and the minimum of \(G\)
over \(\Lambda\) are attained and equal. Every minimizing multiplier
\(\lambda^*\) recovers the weight of every primal maximizer by the formula
\[
          \mathbf z^*=(m-v(\lambda^*))/\Gamma .\tag{10}
\]
The intercept and slacks may have several optimal choices. If
\(\kappa<1/(1-\alpha)\), the dual is infeasible and the primal is unbounded
above. Indeed, all rows together can hold only \(\kappa\) units of mass.
In the primal take \(\mathbf z=0,\ b=t\) and
\(s_i=\max(0,(1-\alpha)t-\min_a g_{ia})\). For all sufficiently large
\(t\) the objective has positive slope \(1-\kappa(1-\alpha)\).

\newpage\subsection*{The duality gap and a primal optimizer}
When total capacity suffices, dividing the required mass equally among rows
and placing each share on one action proves feasibility. The closed,
bounded set \(\Lambda\) is compact, so \(G\) attains a minimum. For any
\(\lambda\in\Lambda\) write \(\mathbf z_\lambda=(m-v(\lambda))/\Gamma\).
With \(\rho_{ia}=g_{ia}+s_i-\langle d_{ia},\mathbf z\rangle-(1-\alpha)b\),
expanding the square gives
\[
 G(\lambda)-P(\mathbf z,b,s)=
 \tfrac{\Gamma}{2}\|\mathbf z-\mathbf z_\lambda\|^2
 +\sum_{i,a}\lambda_{ia}\rho_{ia}
 +\sum_i(\kappa/N-\sum_a\lambda_{ia})s_i .
\]
All three terms are nonnegative at feasible points. It remains to find a
primal point at which all vanish. Fix a minimizer \(\lambda^*\) and set
\(c_{ia}=g_{ia}-\langle d_{ia},\mathbf z_{\lambda^*}\rangle\).
Differentiating \(G\) from \(\lambda^*\) along each segment in \(\Lambda\)
shows that \(\lambda^*\) minimizes the linear cost \(\sum c_{ia}\lambda_{ia}\).

Put \(r_i=\min_a c_{ia}\). To minimize the linear cost, use a cheapest
action in each occupied row. Sort rows by \(r_i\), and let \(u\) be the
cost of the first row at which cumulative capacity reaches \(1/(1-\alpha)\).
Fill all rows cheaper than \(u\), distribute the remaining mass across rows
tied at \(u\), and leave dearer rows empty. This allocation attains the
lower bound
\[
 \sum_{i,a}c_{ia}\lambda_{ia}\ \ge\
       \frac{u}{1-\alpha}-\frac{\kappa}{N}\sum_i(u-r_i)_+ ,
 \qquad (t)_+=\max(t,0).
\]
The bound follows by subtracting \(u\) from each row price and applying the
cap to each negative coefficient. Thus \(\lambda^*\) also attains it.
Set \(b=u/(1-\alpha)\) and \(s_i=(u-r_i)_+\). These choices are feasible
and the two sums in the gap total zero. This proves attainment, and the
squared norm in the gap proves uniqueness of the weight.

\lead{Converting the printed objective.}
In (9) the quadratic is \(F(\lambda)=\frac12\lambda^\top Q\lambda+
R^\top\lambda+S\), where
\(Q_{ia,jc}=\langle d_{ia},d_{jc}\rangle\),
\(R_{ia}=\Gamma g_{ia}-\langle m,d_{ia}\rangle\), and \(S=-N^2\|m\|^2\).
Direct expansion gives
\[
 G(\lambda)=\bigl(F(\lambda)-S+\tfrac12\|m\|^2\bigr)/\Gamma .
\]
Both minimizations give the same multipliers. For one self-looping state
and action, \(\Phi=0,\ \nu=\pi=1,\ g=1,\ \alpha=1/2,\ \kappa=4,\ \Gamma=2\),
the sole feasible multiplier is \(2\). The printed minimum is \(4\), while the
primal maximum and normalized minimum are both \(2\). With \(g=0,\kappa=1\)
the same model realizes the unbounded case.

At the threshold, take two states \(x,y\), one draw at \(x\), and one action
moving to \(y\). The respective costs are \(1,0\), the scalar features \(1,0\),
both laws are uniform, and \(\alpha=1/2,\ \Gamma=\kappa=2\).
The only multiplier is \(2\), giving \(\mathbf z^*=-1/2\) and \(G=9/4\).
The pair \(b=3,\ s_1=0\) attains this value, as does every
\(b=3+t,\ s_1=t/2\) for \(t\ge0\).

\partsection{Kernel evaluation}{Kernel evaluation and greedy actions}
\label{sec:kernel}
Let \(K(x,y)=\langle\Phi(x),\Phi(y)\rangle\). The dual and the policy can
be evaluated with kernel values alone. In the following sums \(x\) may be
any state, including one absent from the sample. For a state and action put
\[
 \mathcal N_{x,a}=\{x\}\cup\{y:p(x,y,a)>0\},\qquad
 q_{x,y,a}=\mathbf1_{\{x=y\}}-\alpha p(x,y,a).
\]
Since \(q_{x,y,a}=0\) outside the neighborhood, the vector \(d_{ia}\) equals
\(\sum_yq_{X_i,y,a}\Phi(y)\). Substitution gives the kernel coefficients
\[
 \begin{aligned}
 Q_{ia,jc}&=\sum_{y,u}q_{X_i,y,a}q_{X_j,u,c}K(y,u),\\
 R_{ia}&=\Gamma g_{ia}-N^{-1}\sum_j w_j\sum_y q_{X_i,y,a}K(X_j,y).
 \end{aligned}\tag{12--13}
\]
The matrix is positive semidefinite because \(\lambda^\top Q\lambda=
\|v(\lambda)\|^2\). The constant \(S\) can be dropped when finding a
minimizer. At such a minimizer \(\lambda^*\), taking the inner product of
(10) with \(\Phi(x)\) gives the prediction
\[
 \widehat J(x)=\Gamma^{-1}\left[
     N^{-1}\sum_j w_jK(x,X_j)-
     \sum_{i,a}\lambda^*_{ia}\left(K(x,X_i)-\alpha\sum_y p(X_i,y,a)K(x,y)\right)
       \right]. \tag{14}
\]
Every optimal primal fit is \(\widehat J+b^*\). Adding \(b^*\) to a value
function adds \(\alpha b^*\) to every action score at a given state.
Hence the two functions have identical sets of greedy actions, ties included.

\lead{Why the current state belongs in the neighborhood.}
Take states \(s,t\) with features \(1,2\). Action \(A\) always moves to
\(s\), action \(B\) always moves to \(t\). Let \(g_{s,A}=5/4,\ g_{s,B}=1\). At \(t\) both actions cost zero. Use one draw at \(s\), uniform \(\pi\),
\(\nu=(3/4,1/4)\), \(\alpha=1/2,\ \kappa=2,\ \Gamma=1\). The two feature
differences are \(1/2\) and \(0\), and \(m=3/2\). For action \(B\),
omitting the current state replaces \(d_{sB}=1-(1/2)2=0\) by \(-1\).
The diagonal entry of \(Q\) changes from \(0\) to \(1\) and \(R_{sB}\)
from \(1\) to \(5/2\). The formula in (12--13) uses the neighborhood of
\S2.3. The description in \S2.5 omits the current state in this example.

Using the full neighborhood gives \(Q=\mathrm{diag}(1/4,0)\) and
\(R=(1/2,1)\). The minimizer is \((2,0)\), so the recovered weight is
\(1/2\) and \(\widehat J=(1/2,1)\). With optimal intercept \(b^*=2\), the
fitted values are \((5/2,3)\). Action costs computed from either prediction
tie at \(s\) and strictly favor \(A\) at \(t\).


\partsection{Regularization}{What changes when a penalty becomes a radius?}
\label{sec:regular}
Consider the population version of the fit. In this section each row is an
actual state, \(\pi\) is a probability distribution, and
\(d_{xa}=\Phi(x)-\alpha P_a\Phi(x)\). The common feasible set is
\[
  s_x\ge0,\qquad \langle d_{xa},\mathbf z\rangle+(1-\alpha)b\le g_{x,a}+s_x
  \quad\text{for every state and action}.
\]
Writing \(m=\sum_x\nu_x\Phi(x)\), let
\(L(\mathbf z,b,s)=\langle m,\mathbf z\rangle+b-\kappa\sum_x\pi_xs_x\).
The two population programs are
\[
\begin{aligned}
 (5)&\quad \max\ L(\mathbf z,b,s)-\tfrac{\Gamma}{2}\|\mathbf z\|^2,\\
 (15)&\quad \max\ L(\mathbf z,b,s)
                  \quad\text{with }\|\mathbf z\|\le C,\ |b|\le B .
\end{aligned}
\]
This distinction matters when comparing objective values. It also matters
at the exact penalty threshold, where an intercept can vary without changing
the penalized objective.

\lead{Penalty and radius.}
For \(\Gamma>0\), (5) attains a finite maximum exactly when
\(\kappa(1-\alpha)\ge1\). All its optimizers have one weight \(\mathbf z^*\).
Choose \(C=\|\mathbf z^*\|\), and choose \(B\) large enough to contain the
intercept of one optimizer. The optimizers of (15) are then precisely the
optimizers of (5) whose intercepts satisfy the box constraint, and
\[
             \max(15)=\max(5)+\tfrac{\Gamma}{2}\|\mathbf z^*\|^2 .
\]
For \(\kappa(1-\alpha)>1\), some finite box contains \emph{all} the
optimizers of (5). At equality every finite box omits some of them.

\lead{Existence and the intercept.}
For fixed weight, the least slack is
\((1-\alpha)(b-u_x)_+\), where
\[
 u_x=\min_a\{g_{x,a}-\langle d_{xa},\mathbf z\rangle\}/(1-\alpha).
\]
The reduced objective as a function of \(b\) has slope one below every
\(u_x\) and slope \(1-\kappa(1-\alpha)\) above every \(u_x\).
Below the penalty threshold the upper tail is unbounded. At or above it,
clipping \(b\) to the interval between the extreme breakpoints preserves
or improves the value. All breakpoints lie in \([-U,U]\), where
\[
 U=\frac{\|g\|_\infty+\max_{x,a}\|d_{xa}\|\,\|\mathbf z\|}{1-\alpha}.
\]
Thus the reduced penalized objective is bounded above by
\(\|m\|\,\|\mathbf z\|+U-\Gamma\|\mathbf z\|^2/2\), which tends to
\(-\infty\) with the norm. Projecting \(\mathbf z\) onto the span of \(m\)
and all \(d_{xa}\) preserves the linear terms and can only reduce the
penalty. A compact weight ball in this finite span and a compact intercept
interval therefore suffice. The reduced objective is continuous there.

\newpage\subsection*{The optimizer set and the value shift}
The feasible set is convex. If two maximizers had different weights,
their midpoint would improve the penalized objective by strict concavity
of the negative squared norm. For any feasible triple with
\(\|\mathbf z\|\le C=\|\mathbf z^*\|\), optimality in (5) now gives
\[
 L(\mathbf z,b,s)\le\max(5)+\tfrac{\Gamma}{2}\|\mathbf z\|^2
                    \le\max(5)+\tfrac{\Gamma}{2}C^2 .
\]
An optimizer of (5) inside the box attains both inequalities. Any
maximizer of (15) must therefore attain both too. In the strict penalty
case the slopes just computed force every optimal intercept into
\([-U,U]\) at \(\mathbf z^*\). At equality increasing \(b\) by \(t\ge0\)
and every slack by \((1-\alpha)t\) preserves feasibility and value.

\lead{Example with unequal values.}
Two absorbing states have respective scalar features and costs \(-1,+1\),
one action, uniform \(\nu=\pi\), and \(\alpha=1/2,\ \kappa=6,\ \Gamma=1\).
The triple \(z=1,\ b=-1,\ s=0\) is feasible. The negative-feature row
requires \(b\le z-2+2s_-\). Thus \(L=b-3(s_-+s_+)\le z-2\).
Subtracting \(z^2/2\) gives at most \(-3/2-(z-1)^2/2\).
The displayed triple attains \(-3/2\) in (5) and \(-1\) in (15) with
\(C=1,\ B\ge1\). Every penalized optimum has \(z=1\), so any common
optimizer gives a value difference of \(1/2\), whatever box contains it.

\lead{Full-rank features and the best error (16).}
If the state feature vectors span a space of dimension \(|\mathcal S|\),
they are linearly independent. Their Gram matrix
\(\mathsf K_{xy}=\langle\Phi(x),\Phi(y)\rangle\) is invertible because
\(c^\top\mathsf Kc=\|\sum_x c_x\Phi(x)\|^2>0\) for nonzero \(c\).
For any target \(J\) take \(c=\mathsf K^{-1}J\) and
\(\mathbf z=\sum_x c_x\Phi(x)\). Then \(J_{\mathbf z,0}=J\).
In particular, once \(C\) contains this weight for \(J^*\), the
best sup-norm error is zero for every \(B\ge0\). For general features the
best error is nonincreasing as either radius grows, by inclusion of the
feasible classes. The matching radius for a particular \(\Gamma\) is fixed
by its optimizer. Exact representation uses a radius large enough for the
interpolating weight.


\partsection{Sampling and the intercept}{The idealized sampled fit and its intercept}
\label{sec:offset}
For the approximation bounds, \(\nu\) is an initial probability distribution.
Equation (6) estimates the feature term by importance weighting and charges
a quadratic penalty. The program analyzed here uses the exact relevance
mean \(\nu^\top J\), a prescribed box for weight and intercept, and
\(\kappa=2/(1-\alpha)\). Its guarantees concern fitted value error.
The allowed values are \(\J=\{J_{\mathbf z,b}:\|\mathbf z\|\le C,\ |b|\le B\}\),
with \(C,B\ge0\). For any realized sample let
\(P_Nf=N^{-1}\sum_i f(X_i)\). Define
\(r(J)=(J-TJ)_+\) and \(r_\mu(J)=(J-T_\mu J)_+\) pointwise.
Since \(TJ\le T_\mu J\), \(0\le r_\mu(J)\le r(J)\).
The idealized program (18) maximizes
\[
  \nu^\top J-\frac{2}{(1-\alpha)N}\sum_i s_i,\qquad
  J\in\J,\quad s_i\ge0,\quad J(X_i)\le TJ(X_i)+s_i .
\]
Its least feasible slack is \(r(J)(X_i)\). The objective decreases
strictly with each slack, so the optimum for each \(J\) fills in exactly
these values. The reduced objective is
\[
                \nu^\top J-\frac{2}{1-\alpha}P_Nr(J). \tag{18}
\]
It has a maximizer on every draw. Projection onto the finite span of the
state features preserves the fitted vector and weakly decreases the weight
norm. Hence \(\J\) is a continuous image of a compact ball times \([-B,B]\). The reduced objective is continuous on that image.

\lead{Exactly which intercepts maximize (18)?}
Hold the weight fixed and put \(V(x)=\langle\Phi(x),\mathbf z\rangle\).
Define the sample breakpoints \(u_i=(TV(X_i)-V(X_i))/(1-\alpha)\).
The identity \(T(V+b)=TV+\alpha b\) gives
\(r(V+b)(X_i)=(1-\alpha)(b-u_i)_+\). Since
\(2(b-u_i)_+=b-u_i+|b-u_i|\), (18) becomes
\[
          \nu^\top V+\frac1N\sum_i u_i-\frac1N\sum_i |b-u_i| .
\]
Thus the optimal free intercepts are exactly the sample medians. In
increasing order they form the interval between entries
\(\lfloor(N+1)/2\rfloor\) and \(\lceil(N+1)/2\rceil\), counting from one.
Between successive distinct entries the slope of the sum of distances
is the number of entries to the left minus the number to the right.
At a tied entry use the two one-sided slopes. This proves the interval
description for odd and even samples, including repeated states.

\newpage\subsection*{How large must the intercept box be?}
Write \(K=\max_x\|\Phi(x)\|\), and \(\|g\|_\infty=\max_{x,a}|g_{x,a}|\).
Every discounted feature difference has norm at most \((1+\alpha)K\), so every breakpoint has absolute value at most
\[
                  B_0=\frac{\|g\|_\infty+(1+\alpha)CK}{1-\alpha}.
\]
Every free maximizer is inside the box when \(B\ge B_0\). Conversely, one
can maximize over the free intercept at any feasible weight by choosing a
median in that box. Hence the boxed and free problems have identical
maximizers, including their slacks, and for \(B>B_0\) every optimizer lies
strictly inside the box.

For a particular weight there is a more precise test. A box that intersects
the median interval preserves the best value at that weight. It preserves
every optimal intercept precisely when it contains the whole interval.
For instance, with two breakpoints \(-1,3\), any \(b\in[-1,3]\) gives the
same sum of distances, \(4\). A box \([-1,1]\) keeps this optimal value
and removes the optimal choices in \((1,3]\).

\lead{A positive-probability failure of the printed bound.}
For the printed rule \(B>2CK+\|g\|_\infty/(1-\alpha)\), take two states
with features \(-1,+1\), zero cost, and a single action leading from
either state to the positive-feature state. Use \(\alpha=1/2,\ C=1\)
and uniform \(\nu\). One draw at the negative-feature state reduces (18) to
\[
 \max\ b-4s \qquad\text{subject to}\qquad
 -1\le z\le1,\quad s\ge0,\quad b-3z\le2s .
\]
With a free intercept, \(b-4s\le3z-2s\le3\), and \((z,b,s)=(1,3,0)\)
attains \(3\). Now take \(B=5/2\). This clears the printed threshold
of \(2\), yet \(b-4s\le5/2\). The triple \((1,5/2,0)\) attains the
smaller value. The draw occurs with probability \(1/4\) under the ideal
distribution in the next section. The example attains the corrected
radius \(B_0=3\).


\partsection{Occupancy and sampling}{Discounted visits and uniform sampling error}
\label{sec:statistics}
Let \(\nu\) be the starting distribution. For a policy \(\mu\), its
discounted occupancy distribution assigns a state its expected discounted
fraction of visits according to
\[
 \pi_{\mu,\nu}^{\top}
 =(1-\alpha)\sum_{t=0}^{\infty}\alpha^t\nu^\top P_\mu^t
 =(1-\alpha)\nu^\top(I-\alpha P_\mu)^{-1}. \tag{17}
\]
Each \(\nu^\top P_\mu^t\) is a probability row and the coefficients sum
to one. The matrix series converges entrywise. Multiplying the partial
sum by \(I-\alpha P_\mu\) gives \(I-\alpha^{n+1}P_\mu^{n+1}\to I\).
Thus (17) is also the unique solution of
\(\pi_{\mu,\nu}^\top=(1-\alpha)\nu^\top+\alpha\pi_{\mu,\nu}^\top P_\mu\).
Subtracting \(J_\mu=T_\mu J_\mu\) from the one-step equation yields
\[
 \nu^\top(J-J_\mu)=\pi_{\mu,\nu}^\top(J-T_\mu J)/(1-\alpha).
\]
In the example of the preceding section, only time zero can contribute a
visit to the negative-feature state, giving mass \((1-\alpha)\nu_-=1/4\).

\lead{Theorem 3 (uniform sampling for one policy).}
Fix a policy \(\mu\) and any law \(\pi\) on the states. Draw
\(X_1,\ldots,X_N\) independently from \(\pi\). Write \(Pf=\pi^\top f\).
For \(0<\delta<1,\ \epsilon>0\), put
\[
 \begin{aligned}
 \Cbar&=4CK(1+\alpha)+4B(1-\alpha)+2\|g\|_\infty,\\
 \Xi&=\Cbar\bigl(1+\sqrt{2\log(1/\delta)}\bigr).
 \end{aligned}
\]
Here and throughout the approximation bounds \(K=\max_x\|\Phi(x)\|\).
If \(N\ge1\) and \(N\ge\Xi^2/\epsilon^2\), then
\[
 \Prob\left\{\sup_{J\in\J}\bigl(P r_\mu(J)-P_Nr_\mu(J)\bigr)
                   \ge\epsilon\right\}\le\delta .\tag{T3}
\]
The function can be fitted on the same draws. Only the policy and the
sampling law are fixed in advance.

\lead{Pointwise bounds (Lemma 3).}
Let \(d_\mu(x)=\Phi(x)-\alpha P_\mu\Phi(x)\). Then
\[
 r_\mu(J_{\mathbf z,b})(x)
   =(\langle d_\mu(x),\mathbf z\rangle+(1-\alpha)b-g_\mu(x))_+ .
\]
Since \(\|d_\mu(x)\|\le(1+\alpha)K\), both \(r_\mu(J)\) and its
optimal-action counterpart \(r(J)\) lie in \([0,\Cbar/2]\).
The latter is the pointwise maximum of the action residuals.

\lead{A check on the signed identity.}
Start at \(s\) in a two-state model with \(\alpha=1/2\). Under the given
policy, \(s\) costs \(1\) and moves uniformly to \(s,t\), while \(t\) costs
\(2\) and stays at \(t\). Balance gives \(\pi_{\mu,\nu}=(2/3,1/3)\);
policy evaluation gives \(J_\mu=(8/3,4)\). For \(J=(3,2)\), the signed
residuals are \(3/4,-1\). Their occupancy mean is \(1/6\). Dividing by
\(1-\alpha\) recovers the initial error \(3-8/3=1/3\).

\newpage\subsection*{Bounding the complexity of the residual class}
For a nonempty class \(\mathcal F\) write
\[
 \widehat R_N(\mathcal F)=
   \E_\sigma\sup_{f\in\mathcal F}\frac{2}{N}\sum_i\sigma_i f(X_i),
 \qquad R_N(\mathcal F)=\E_X\widehat R_N(\mathcal F).
\]
The signs are independent, uniform on \(\{-1,1\}\), and independent of the
sample. This is the manuscript's normalization.

\lead{Lemma 2 and Corollary 1.}
For \(x\mapsto\langle d_\mu(x),\mathbf z\rangle\), the empirical complexity
is at most \(2CK(1+\alpha)/\sqrt N\). Conditional on the sample,
Cauchy--Schwarz bounds the signed supremum by
\((2C/N)\|\sum_i\sigma_i d_\mu(X_i)\|\). The expected squared norm of this
sum is \(\sum_i\|d_\mu(X_i)\|^2\), since distinct signs have zero product
in expectation. The intercept contributes at most
\(2(1-\alpha)B/\sqrt N\) by the same argument for scalar signs.

Subtracting the state cost and taking the positive part applies a
1-Lipschitz function to each coordinate. To check contraction for this
one-sided supremum, condition on every sign but one. For candidate \(j\),
let \(u_j\) be its score at that coordinate and \(a_j\) the signed sum of
the other scores. Twice the conditional expectation for that coordinate is
\(\sup_{j,k}(a_j+a_k+\varphi(u_j)-\varphi(u_k))\).
Replacing the last difference by \(|u_j-u_k|\) bounds it above.
Exchanging \(j,k\) makes the bound equal to the expression for the
identity function. Repeat coordinate by coordinate. Thus for
\(\mathcal F_\mu=\{r_\mu(J):J\in\J\}\),
\[
 \widehat R_N(\mathcal F_\mu)
 \le 2[CK(1+\alpha)+B(1-\alpha)]/\sqrt N\le\Cbar/\sqrt N .
\]
Taking expectations proves Corollary~1.

\lead{Lemma 1 (one-sided concentration).}
For independent draws from the same law \(P\) on a finite or countable discrete state space,
suppose \(\mathcal F\) is nonempty and \(|f(x)|\le B_f\) for every \(f,x\),
with \(B_f>0\). For \(Z=\sup_{f\in\mathcal F}(Pf-P_Nf)\),
\[
 \Prob\{Z\ge R_N(\mathcal F)+B_f\sqrt{2\log(1/\delta)/N}\}\le\delta .
\]
To reduce an uncountable class to a countable one, use the
pointwise topology and choose one function from each nonempty basic open set
given by rational intervals on finitely many states. Empirical means
are continuous, and bounded convergence gives continuity of the population
mean on this uniformly bounded class. Hence the countable dense subclass
has the same suprema.

\newpage\subsection*{From symmetrization to a probability bound}
An independent sample with empirical mean \(P_N'\) gives
\(\E Z\le\E\sup_f(P_N'f-P_Nf)\), by bringing the supremum inside the
expectation. Swapping each pair of observations independently preserves
its law and inserts a uniform sign into each difference. Splitting the
signed sum at the supremum and using symmetry gives \(\E Z\le R_N(\mathcal F)\).
Changing one observation changes \(Z\) by at most \(2B_f/N\).

For completeness, if changing coordinate \(i\) of an independent sample
changes a statistic by at most \(c_i\), the successive conditional-mean
increments are centered and have conditional ranges of width at most \(c_i\).
For such an increment the second derivative of the log moment-generating
function is a tilted variance, at most \(c_i^2/4\). Integrating twice and
iterating conditional expectations bounds the centered exponential moment
by \(\exp(\lambda^2\sum_i c_i^2/8)\). Markov's inequality optimized at
\(\lambda=4t/\sum_i c_i^2\) gives
\(\Prob(Z-\E Z\ge t)\le\exp(-2t^2/\sum_i c_i^2)\) for \(t>0\) and a
positive denominator. With \(c_i=2B_f/N\) this proves Lemma~1. If \(B_f=0\)
and \(\mathcal F=\{0\}\), its printed closed event has probability one.

\lead{Finishing Theorem 3 and the fixed-function part of Lemma 3.}
For \(\Cbar>0\), apply Lemma~1 with \(B_f=\Cbar/2\) to \(\mathcal F_\mu\).
The threshold in Lemma~1 is at most \(\Xi/\sqrt N\le\epsilon\), so the event
in (T3), including equality, has probability at most \(\delta\).
For a single fixed \(f\) from either residual class,
\(P_N f\) changes by at most \(\Cbar/(2N)\) in each coordinate.
The same tail bound now gives
\[
 \Prob(P_Nf-Pf\ge\epsilon)\le
   \exp(-8N\epsilon^2/\Cbar^2)\le\delta^4 .
\]
If \(\Cbar=0\), all the residuals vanish. Since \(\epsilon>0\), both failure
events are empty. The sample bounds depend on the radius and on feature
norms, without a factor for the dimension of the ambient Hilbert space.


\partsection{Approximation guarantees}{The approximation error of the fitted value}
\label{sec:guarantees}
Choose an optimal policy \(\mu^*\) from Section~\ref{sec:bellman}.
Throughout this section \(X_i\) are independent draws from
\(\pi=\occ\), \(Pf=\pi^\top f\), and \(\widehat J\) is any optimizer of
(18). Use the sample threshold and confidence parameters of Theorem~3.
Measure error by \(\|U\|_{1,\nu}=\sum_x\nu_x|U(x)|\), and let
\(\|U\|_\infty=\max_x|U(x)|\).

\lead{Theorem 1 (uniform comparator error).}
On an event of probability at least \(1-\delta\), every optimizer satisfies
\[
 \|J^*-\widehat J\|_{1,\nu}\le
 \frac{3+\alpha}{1-\alpha}\min_{V\in\J}\|J^*-V\|_\infty
            +\frac{2\epsilon}{1-\alpha}. \tag{19$'$}
\]
In particular this implies the printed (19), which has confidence
\(1-\delta-\delta^4\) and additive term \(4\epsilon/(1-\alpha)\).

\lead{A common inequality for both theorems.}
Let \(D=(I-\alpha P_{\mu^*})^{-1}\), which has nonnegative entries by (17).
For any \(J\), \(J-J^*=D(J-T_{\mu^*}J)\), hence
\((J-J^*)_+\le D r_{\mu^*}(J)\). Using
\(|u|=-u+2u_+\) and \(\nu^\top D=\pi^\top/(1-\alpha)\) gives
\[
 \|J^*-J\|_{1,\nu}\le\nu^\top(J^*-J)+\tfrac{2}{1-\alpha}P r_{\mu^*}(J).
\]
On the good event of (T3), the last mean for the fitted value is at most
\(P_Nr(\widehat J)+\epsilon\). Optimality in (18) compares its reduced
objective with that of every \(V\in\J\), giving
\[
 \|J^*-\widehat J\|_{1,\nu}
 \le\nu^\top(J^*-V)+\tfrac{2}{1-\alpha}\bigl(P_Nr(V)+\epsilon\bigr).
 \tag{*}\label{eq:oracle}
\]
The comparator needs slack \((V-TV)_+\). For a one-state zero-cost
self-loop, \(\alpha=1/2,\ V=1\), this is \(1/2\). Reversing the sign as in
Appendix~B.2 gives zero and violates the constraint.

\lead{Proof of Theorem 1.}
Let \(d=\|J^*-V\|_\infty\). Contraction of \(T\) and \(TJ^*=J^*\) imply
\(r(V)(x)\le(1+\alpha)d\) at every state. Also \(\nu^\top(J^*-V)\le d\).
Substitution in \eqref{eq:oracle} gives (19$'$) by minimizing over the
compact set of fits. This step requires only the event for the fitted
function.

\newpage\subsection*{A sharper comparison when errors vary across states}
For a positive state weight \(\psi(x)\ge1\), put
\[
 \|U\|_{\infty,1/\psi}=\max_x |U(x)|/\psi(x),\qquad
 \beta(\psi)=\max_{x,a}\frac{P_a\psi(x)}{\psi(x)} .
\]
A larger \(\psi(x)\) allows greater error at \(x\). The two means of
\(\psi\) in the guarantee charge for this under initial and discounted visits.

\lead{Theorem 2 (weighted comparison).}
With probability at least \(1-\delta-\delta^4\), every optimizer satisfies
\[
\begin{aligned}
 \|J^*-\widehat J\|_{1,\nu}
 &\le \inf_{\substack{V\in\J\\\psi\ge1}}
  \left(\nu^\top\psi+
    \frac{2(\pi^\top\psi)(1+\alpha\beta(\psi))}{1-\alpha}\right)
            \|J^*-V\|_{\infty,1/\psi}
       +\frac{4\epsilon}{1-\alpha}.
\end{aligned}\tag{25}
\]
For \(\psi\equiv1\) its coefficient is \((3+\alpha)/(1-\alpha)\).
For the nontrivial confidence bounds we use \(0<\delta<1\). The
printed probability lower bounds for Theorems~1 and~2 are vacuous if
\(\delta\ge1\).

\lead{Proof.}
Define \(L(V)=\nu^\top(J^*-V)+2Pr(V)/(1-\alpha)\). Choose a minimizer
\(V_0\) of this continuous functional on \(\J\) \emph{before} the draws.
The fixed-function part of Lemma~3 then gives
\(P_Nr(V_0)\le Pr(V_0)+\epsilon\) outside an event of probability
at most \(\delta^4\). Intersect this good event with that of
\eqref{eq:oracle}. For every optimizer the error is at most
\(L(V_0)+4\epsilon/(1-\alpha)\).

For any \(V,\psi\), put \(d=\|J^*-V\|_{\infty,1/\psi}\). For every action,
\[
 |\alpha P_a(J^*-V)(x)|\le\alpha d\,P_a\psi(x)
                      \le\alpha d\,\beta(\psi)\psi(x).
\]
Taking the minimum over actions preserves this last bound on the
difference. Hence \(r(V)(x)\le(1+\alpha\beta(\psi))d\psi(x)\), while
\(J^*(x)-V(x)\le d\psi(x)\). Averaging gives the comparison term in
(25) as an upper bound for \(L(V)\) and hence for \(L(V_0)\).
This holds for every pair \(V,\psi\) on the same event. Taking the
infimum proves the result, even when no weight attains it.

\lead{The infimum can be unattained.}
Two absorbing states have costs \(0,5\), discount \(1/2\), and initial
weights \(9/10,1/10\). Then \(\pi=\nu,\ J^*=(0,10)\). Take
\(C=B=0\), so \(V=0\). For \(\psi=(u,v)\ge1\) the comparison term is
\(7+63u/v\). Its infimum is \(7\), approached with \(u=1,\ v\to\infty\).


\partsection{Change of sampling law}{Approximation under a different sampling law}
\label{sec:samplinglaw}
Keep the initial distribution \(\nu\) and an optimal policy \(\mu^*\).
Write \(\pi=\occ\) for the ideal law, but draw the data in (18) independently
from an arbitrary probability distribution \(\rho\). Use \(\rho_N f\)
for the observed mean. Theorem~3 still controls the difference between
\(\rho^\top f\) and \(\rho_Nf\). The loss calculation uses \(\pi^\top f\),
so we need to pass between these two \emph{population} means.

\lead{Total variation.}
Use the convention
\(\mathrm{TV}(\pi,\rho)=\tfrac12\sum_x|\pi_x-\rho_x|\).
For \(0\le f\le L\), subtracting \(L/2\) from \(f\) leaves the difference
of means unchanged and gives
\[
 |(\pi-\rho)^\top f|\le L\,\mathrm{TV}(\pi,\rho).
\]
Opposite point masses and the indicator of the first point attain equality.
This inequality also holds for two realized empirical distributions.
Consequently, when a difference of nonnegative slack means is multiplied
by \(2/(1-\alpha)\), its sharp general TV coefficient is
\(2\|s\|_\infty/(1-\alpha)\). The coefficient printed in Appendix~B.3 is
half that under the convention above.

\lead{Approximation guarantees under the observed law.}
Let \(L=\Cbar/2,\ \tau=\mathrm{TV}(\pi,\rho)\), with the same
\(\epsilon,\delta,N\) as in Theorem~3. With probability at least
\(1-\delta\), every optimizer now satisfies
\[
 \|J^*-\widehat J\|_{1,\nu}\le
 \frac{3+\alpha}{1-\alpha}\min_{V\in\J}\|J^*-V\|_\infty
          +\frac{2}{1-\alpha}(\epsilon+L\tau).
\]
Indeed for the fixed-policy fitted residual \(f=r_{\mu^*}(\widehat J)\),
\(\pi^\top f\le\rho^\top f+L\tau\le\rho_Nf+\epsilon+L\tau\).
Insert this estimate in the proof of \eqref{eq:oracle}.
For any comparator \(V\) fixed before sampling, Lemma~3 also gives
\(\rho_Nr(V)\le\rho^\top r(V)+\epsilon\) with failure probability
at most \(\delta^4\). Therefore with confidence \(1-\delta-\delta^4\),
\[
 \|J^*-\widehat J\|_{1,\nu}\le\nu^\top(J^*-V)
 +\frac{2}{1-\alpha}\bigl[\rho^\top r(V)+2\epsilon+L\tau\bigr].
\]
One may take the deterministic minimizer of
\(\nu^\top(J^*-V)+2\rho^\top r(V)/(1-\alpha)\) on \(\J\).
The weighted comparison in (25) then carries over with
\(\rho^\top\psi\) in place of \(\pi^\top\psi\) and this extra TV term.

\lead{What happens to two empirical distributions?}
For a fixed \(N\), convergence also holds in probability under an explicit
coupling. Put \(m_x=\min(\pi_x,\rho_x)\). Its total mass is \(1-\tau\).
For one pair put mass \(m_x\) at \((x,x)\). When \(\tau>0\), add mass
\((\pi_x-m_x)(\rho_y-m_y)/\tau\) at \((x,y)\). The marginals are
\(\pi,\rho\), and the added mass is off the diagonal. At \(\tau=0\) use only
the diagonal. Draw \(N\) pairs independently. All coordinates agree with
probability \((1-\tau)^N\), so the probability that the two empirical laws
differ is at most \(1-(1-\tau)^N\le N\tau\). With independent sampling, two
single draws from the same fair two-state law differ with probability \(1/2\).

\newpage\subsection*{An observable correction under a ratio bound}
Assume \(\pi_x\le(1+k)\rho_x\) at every state for some \(k\ge0\).
For any \(f\ge0\), \(\pi^\top f\le(1+k)\rho^\top f\).
For \(f=r_{\mu^*}(\widehat J)\), write \(\overline s=N^{-1}\sum_i\widehat s_i\)
for the observed optimal mean slack. It dominates \(\rho_Nf\).
On the same uniform event,
\[
 \pi^\top f\le \rho_N f+(1+k)\epsilon+k\overline s.
\]
The fixed-comparator guarantee above therefore holds with
\(2\epsilon+L\tau\) replaced by \((2+k)\epsilon+k\overline s\), at the
same confidence. To read the two corrections alongside (25), define
\[
 W_\rho=\inf_{\substack{V\in\J\\\psi\ge1}}
 \left(\nu^\top\psi+\frac{2(\rho^\top\psi)(1+\alpha\beta(\psi))}{1-\alpha}\right)
             \|J^*-V\|_{\infty,1/\psi}.
\]
Taking the deterministic population minimizer as in Theorem~2 gives, at
confidence \(1-\delta-\delta^4\), the corresponding weighted bounds
\[
 \|J^*-\widehat J\|_{1,\nu}\le W_\rho+
 \frac{2}{1-\alpha}
 \begin{cases}
  2\epsilon+L\tau & \text{for arbitrary }\rho,\\
  (2+k)\epsilon+k\overline s & \text{when }\pi\le(1+k)\rho .
 \end{cases}
\]

\lead{A counterexample for empirical slacks.}
Appendix~B.3 instead proposes coupling two samples so that
\((\pi_N-\rho_N)^\top\widehat s\le k\rho_N^\top\widehat s\) with high
probability, where \(\widehat s(x)=r(\widehat J)(x)\) is defined on
all states and \(\pi_N,\rho_N\) denote their two empirical distributions.
The following example applies to every coupling, meaning any joint draw
with the prescribed i.i.d. marginal laws.

Take two absorbing states, discount \(1/2\), one action, and costs \(-1,+1\).
Call the negative-cost state \(s\). Put \(\Phi=0,\ C=B=1\). Every fit is a
constant \(b\in[-1,1]\) and its least statewise slack is
\((1+b/2)\mathbf1_{\{s\}}\). Its positive factor cancels from the proposed
comparison even when \(b\) depends on the sample. Let \(\rho_s=1/2\) and
\(\pi_s=\nu_s=1-t\) with \(0\le t\le1\). Absorption gives this ideal occupancy exactly, and
\(\pi\le2\rho\). Thus \(k=1\) suffices.

For odd \(N\), the observed sample visits \(s\) fewer than \(N/2\) times
with probability exactly \(1/2\), by symmetry. The ideal sample consists
entirely of \(s\) with probability at least \(1-Nt\). Under any coupling
the intersection has probability at least \(1/2-Nt\). On it the ideal
indicator mean is one and twice the observed mean is strictly smaller.
For \(t=0\) the failure probability is exactly \(1/2\) at every odd
sample size. For \(t=1/(4N)\) both laws have full support and the failure
probability is at least \(1/4\). The latter is a uniform obstruction over
positive laws. The extra sampling term in the corrected guarantee
accounts for fluctuations even when the observed slacks are small.


\partsection{The active-set algorithm}{Two-coordinate descent and the correct update}
\label{sec:active}
Return to the sampled quadratic \(F\) and its feasible set \(\Lambda\)
in Section~\ref{sec:dual}. Write \(h=Q\lambda+R\) for the gradient at a
feasible multiplier. A coordinate is a draw--action pair \(p=(i,a)\), and
\(e_p\) denotes its unit vector. A transfer \(e_p-e_q\) moves mass from
\(q=(j,c)\) to a distinct \(p\). A positive step is legal precisely when
\(\lambda_q>0\) and either \(i=j\) or row \(i\) is below its cap.

\lead{Proposition 2.}
At every suboptimal feasible multiplier there is a legal pair with
\(h_p<h_q\). All sufficiently small positive steps along \(e_p-e_q\)
strictly improve \(F\). A feasible point is globally optimal exactly when
all legal transfers have nonnegative slope.

\lead{Lemma 4 (exchange).}
For finite vectors with \(\sum_i x_i=0,\ \sum_i x_i y_i<0\), some indices
satisfy \(x_i>0,\ x_j<0,\ y_i<y_j\). The positive and negative parts of
\(x\) carry the same nonzero mass. Dividing by this mass shows that the
weighted average of \(y\) on the positive coordinates is smaller than its
weighted average on the negative ones. At least one such pair exists.

\lead{Proof of Proposition 2.}
For any direction \(d\),
\[
 F(\lambda+td)-F(\lambda)=t h^\top d+\tfrac12 t^2 d^\top Qd.
\]
If \(\lambda'\) is feasible and better, \(h^\top(\lambda'-\lambda)<0\)
by positive semidefiniteness of \(Q\). Suppose every legal transfer has
nonnegative slope. In each row \(i\) every positive coordinate must have
the minimum score \(m_i=\min_a h_{ia}\), since a within-row transfer can
move mass to any action there. Let \(x_i=\sum_a(\lambda'_{ia}-\lambda_{ia})\).
Then \(\sum_i x_i=0\), and
\(\sum_i m_i x_i\le h^\top(\lambda'-\lambda)<0\). Lemma~4 gives a gaining
row \(i\) with smaller \(m_i\) than some losing row \(j\). The gaining row
has room, and the losing row has a positive coordinate of score \(m_j\).
Transferring from it to a cheapest action in \(i\) is legal and has negative
slope, contradicting the assumption. The exact quadratic expansion shows
that all sufficiently small positive steps for this pair decrease \(F\).

\newpage\subsection*{Selecting and updating the pair}
Under the normalization \(\|d\|_\infty\le1\) and sparsity at most two
in (20), the total-mass constraint forces each nonzero direction to have
opposite entries \(t,-t\) with \(0<t\le1\). A negative slope is minimized
at \(t=1\). For a receiver in a row with room, compare the lowest score
in those rows with the largest score among all positive coordinates.
For a receiver in a full row, compare its lowest score with the largest
score among the positive coordinates of that same row. Skip empty groups
and choose the most negative difference across all candidates.
A negative difference automatically picks distinct coordinates. If none
exists, zero solves (20) and Proposition~2 certifies global optimality.

With \(\mathcal B=\{p,q\}\), freeze \(\lambda_{\bar{\mathcal B}}\) and
let \(u=\lambda_{\mathcal B}\). The exact restriction of \(F\), up to a
constant, is
\[
 \tfrac12u^\top Q_{\mathcal B\mathcal B}u+
 u^\top(R_{\mathcal B}+Q_{\mathcal B\bar{\mathcal B}}\lambda_{\bar{\mathcal B}}).
 \tag{22$'$}
\]
Indeed the two symmetric off-diagonal blocks have equal contributions.
The printed (22) doubles the quadratic and cross terms and keeps the same
linear cost. One can also solve along the transfer. Let \(t_{\max}\)
be the donor's mass, further limited by the receiver's available row capacity
for a cross-row move. Put \(c=(e_p-e_q)^\top Q(e_p-e_q)\ge0\). If
\(h_p<h_q\), the minimizing step on \([0,t_{\max}]\) is
\[
 t^*=\begin{cases}
 \min\{t_{\max},(h_q-h_p)/c\}, & c>0,\\
 t_{\max}, & c=0 .
 \end{cases}
\]
This proves strict descent and the stated stopping test.

\lead{A stalled update in (22).}
Take two states \(s,t\), actions \(A,B\) moving to \(s,t\), respectively,
and one draw at \(s\). The actions cost \(0,-1\) at \(s\). Use
\(\Phi(s)=0,\ \Phi(t)=2,\ \alpha=1/2,\ \Gamma=1,\ \kappa=2\), with
uniform relevance and sampling laws. Then \(m=0\), the feature differences
are \(0,-1\), \(Q=\mathrm{diag}(0,1)\), and \(R=(0,-1)\). All feasible
multipliers have the form \((2-u,u)\) for \(0\le u\le2\). Completing
the square shows where the two updates go:
\[
\begin{aligned}
 \text{Restriction of (9)} &\qquad \tfrac12u^2-u
                         =\tfrac12(u-1)^2-\tfrac12 ,\\
 \text{Printed update (22)}&\qquad u^2-u=(u-\tfrac12)^2-\tfrac14 .
\end{aligned}
\]
Starting at \(u=1/2\), the printed update returns \(1/2\) despite the
available decrease. The correct update goes to \(u=1\), lowering the
true objective from \(-3/8\) to \(-1/2\).


\partsection{Corrections and proof sources}{Corrections to the manuscript and formal proofs}
\label{sec:errata}
\lead{The numbered results.}
Theorems~1--3, Proposition~2, Lemmas~2--4, and Corollary~1 hold under the
hypotheses stated in these sections. Proposition~1 needs a penalty
threshold and a normalization of the dual value. Lemma~1 needs a positive
envelope for its closed failure event. The zero-envelope case of Theorem~3
holds directly. The following entries locate the changes a reader would
make when revising the manuscript.

\begin{description}
\item[\S2.1 and (1).] Minimize discounted cost over policies. The Bellman
LP maximizes a pointwise lower bound and uniquely recovers the cost when
all relevance weights are positive. See Section~\ref{sec:bellman}.
\item[Proposition 1 and \S2.5.] Require \(\kappa\ge1/(1-\alpha)\) and
convert the objective value of (9) before equating it to the primal
maximum. Include the current state in the kernel neighborhood.
See Sections~\ref{sec:dual} and~\ref{sec:kernel}.
\item[\S3.1 and (15).] Match \(C\) to the optimal weight norm, add the
quadratic penalty back when comparing values, and distinguish full equality
of optimizer sets from equality inside the intercept box at the threshold.
Full-rank interpolation remains valid. See Section~\ref{sec:regular}.
\item[\S3.2 and (18).] The optimal free intercept is a sample median.
The corrected uniform intercept radius can exceed the printed bound. Theorems~1--3
use the specified bounded class and retain their guarantees for it.
See Sections~\ref{sec:offset}--\ref{sec:guarantees}.
\item[Appendices B.1, B.2 and C.] Use the positive fixed-policy residual
class and its envelope \(\Cbar/2\) in Theorem~3. The final display on
manuscript p.~30 uses the linear class and a different bound symbol. Use \((V-TV)_+\) for the comparator slack, and choose a
population comparator before sampling for the weighted bound.
See Sections~\ref{sec:statistics} and~\ref{sec:guarantees}.
\item[Appendix B.3.] With half-\(\ell_1\) total variation the printed
coefficient is too small. Population likelihood-ratio domination alone
does not deliver the proposed empirical inequality. The replacement
bounds quantify the extra sampling error. See Section~\ref{sec:samplinglaw}.
\item[\S4.1--4.2.] Skip unavailable pair candidates. Use the exact
restricted quadratic (22$'$) for the update. Equation (22) can stall at a
suboptimal point. See Section~\ref{sec:active}.
\end{description}
\newpage\subsection*{Sources and reproducibility}
The manuscript is N.~Bhat, V.~F.~Farias, C.~C.~Moallemi and A.~T.~Zheng,
\href{https://moallemi.com/ciamac/papers/kernel-adp-2012.pdf}{\emph{Non-parametric Approximate Dynamic Programming via the Kernel Method}},
version of August 29, 2022. The
\href{https://github.com/nikhilpb/kadp}{proof repository} contains the
Lean files, the source of this document and the independently readable
proof notes. State and action spaces are finite in the dynamic programming
results. Lemma~1 also covers countably infinite discrete state spaces.

\begin{center}\small
\begin{tabular}{@{}>{\raggedright\arraybackslash}p{.42\linewidth}>{\raggedright\arraybackslash}p{.52\linewidth}@{}}
\toprule
Result & Lean sources\\ \midrule
Bellman LP and policy convention & \leanfile{BellmanLP}{Bellman LP and cost example}\\
Proposition 1 and kernel policy &
 \leanfile{PrimalAttainment}{Primal attainment},
 \leanfile{Duality}{duality gap},
 \leanfile{KernelPolicy}{kernel evaluation}\\
Population regularization and full rank &
 \leanfile{Regularization}{Optimal sets, values and interpolation}\\
Intercept in (18) & \leanfile{Offset}{Median criterion and counterexample}\\
(17), Lemmas 1--3, Corollary 1 and Theorem 3 &
 \leanfile{Occupancy}{Occupancy},
 \leanfile{OneSidedConcentration}{one-sided tail},
 \leanfile{ResidualComplexity}{residual complexity},
 \leanfile{UniformDeviation}{fixed comparator},
 \leanfile{Theorem3}{Theorem 3}\\
Theorems 1 and 2 &
 \leanfile{Theorem1}{Unweighted bound},
 \leanfile{Theorem2}{weighted bound},
 \leanfile{Examples/Theorem2}{unattained infimum}\\
Change of sampling law &
 \leanfile{SamplingLaw}{Transport bounds and ratio counterexample},
 \leanfile{WeightedSampling}{weighted guarantees}, \leanfile{SamplingCoupling}{empirical coupling}\\
Lemma 4 and Proposition 2 &
 \leanfile{Exchange}{Exchange},
 \leanfile{ActiveSet}{two-coordinate descent},
 \leanfile{PairSelection}{pair search},
 \leanfile{Examples/Proposition2}{stalling example}\\
\bottomrule
\end{tabular}
\end{center}

The Lean package uses \href{https://github.com/leanprover-community/mathlib4}{mathlib}
for finite sums, Hilbert spaces and probability,
\href{https://github.com/LeanEconomics/DynamicProgramming/tree/9d4783f3fb282444f8f3a91c4e6439293dd4bacc}{LeanEconomics}
for Bellman operators, policy values and discounted dynamics, and
\href{https://github.com/auto-res/lean-rademacher/tree/db9f18714b0b9934553c37db1b41250462dc7faa}{FoML}
for concentration and Rademacher complexity. The repository pins the
library revisions and Lean toolchain. Running \texttt{lake build} checks
the declarations. The repository's \href{https://github.com/nikhilpb/kadp/blob/main/README.md}{reproduction instructions}
also give the kernel replay and axiom checks.
\end{document}
